% GENERATED, edit the French unit or the translation table

\mxunit{stmg-c06}{Conditional probability}
\mxskills{model with a tree or a table, calculate, reason about independence, interpret}
\section{Reference course}
For two events $A$ and $B$ with $P(A)>0$, the probability of $B$ given $A$ is
\[
P_A(B)=P(B\mid A)=\frac{P(A\cap B)}{P(A)}.
\]

\begin{center}
\begin{tikzpicture}[grow=right,level distance=34mm,sibling distance=13mm,every node/.style={font=\small},edge from parent/.style={draw,-{Stealth}}]
\node {$\Omega$}
 child {node {$A$}
   child {node {$B$} edge from parent node[above] {$P_A(B)$}}
   child {node {$\overline B$} edge from parent node[below] {$1-P_A(B)$}}
   edge from parent node[above] {$P(A)$}}
 child {node {$\overline A$}
   child {node {$B$} edge from parent node[above] {$P_{\overline A}(B)$}}
   child {node {$\overline B$} edge from parent node[below] {$1-P_{\overline A}(B)$}}
   edge from parent node[below] {$1-P(A)$}};
\end{tikzpicture}
\end{center}

\begin{keypoint}
In a weighted tree, the probabilities leaving the same node sum to $1$. We multiply the probabilities along a path: $P(A\cap B)=P(A)P_A(B)$. We add the probabilities of mutually exclusive paths leading to the event we are looking for.
\end{keypoint}

If $A_1,\ldots,A_n$ form a partition of the sample space, then
\[
P(B)=P(A_1)P_{A_1}(B)+\cdots+P(A_n)P_{A_n}(B).
\]
Two events $A$ and $B$ with nonzero probabilities are independent if $P(A\cap B)=P(A)P(B)$, or equivalently if $P_A(B)=P(B)$.

\begin{pitfall}
Do not confuse $P_A(B)$ and $P_B(A)$. The condition comes after the word ``given'' and is written as a subscript in this notation.
\end{pitfall}

\begin{example}
$40\,\%$ of customers buy online. Among them, $12\,\%$ return their order; among the others, $5\,\%$ return it. Then $P(R)=0.40\times0.12+0.60\times0.05=0.078$. The overall return rate is therefore $7.8\,\%$.
\end{example}

\mxsession{39}{\mxdate*{2027-03-09}}{B: 2 h -- C: 1 h}{Conditioning}{Conditional probability}
\begin{objectives}\begin{itemize}\item Define a conditional probability.\item Move from counts to probabilities.\item Interpret.\end{itemize}\end{objectives}
\begin{fluency}[6 to 8 minutes, no calculator]\begin{enumerate}\item Calculate $20\,\%$ of $140$. \mxblank{$28$}\item What is the sum of the branches leaving a node? \mxblank{$1$}\item Calculate $0.7\times0.4$. \mxblank{$0.28$}\item If $P(A)=0.37$, calculate $P(\overline A)$. \mxblank{$0.63$}\item Simplify $\dfrac{18}{30}$. \mxblank{$\dfrac35$}\end{enumerate}\end{fluency}
\begin{prerequisites}Reread the last written record and be able to carry out the basic calculations used today. Before the investigation, have the keywords restated: define a conditional probability and move from counts to probabilities.\end{prerequisites}

\begin{activity}
In a survey, $180$ customers out of $600$ use a premium card and $72$ of these $180$ customers buy every week. Calculate $P(P)$, $P(P\cap H)$ and $P_P(H)$, then interpret.

\end{activity}
\begin{solution}$P(P)=0.30$, $P(P\cap H)=0.12$ and $P_P(H)=\dfrac{72}{180}=0.40$. Have students verbalise the piece of data that triggers the choice of method.\end{solution}

\begin{summary}For $P(A)>0$, $P_A(B)=\dfrac{P(A\cap B)}{P(A)}$. This is the probability of $B$ within the population restricted to $A$.\end{summary}

\begin{example}
Take the data of the investigation and write the solution in order: definition of the objects, property used, exact or approximate calculation, then a sentence giving the answer. The model answer appears in the teacher edition.
\end{example}
\begin{solution}$P(P)=0.30$, $P(P\cap H)=0.12$ and $P_P(H)=\dfrac{72}{180}=0.40$.\end{solution}

\begin{exercise}[essential]
\begin{enumerate}\item We know $P(A)=0.35$, $P_A(B)=0.6$ and $P_{\overline A}(B)=0.8$. Construct the tree.\item Calculate $P(A\cap B)$ then $P(B)$.\item Calculate $P_B(A)$ and test independence.\end{enumerate}
\end{exercise}
\begin{solution}$P(A\cap B)=0.35\times0.6$. Then $P(B)=P(A)P_A(B)+(1-P(A))P_{\overline A}(B)$ and $P_B(A)=\dfrac{P(A\cap B)}{P(B)}$. Compare with $P(A)$.\end{solution}

\mxfurther{Calao, course pp. 132--137; exercises p. 138, nos.~1--5}
\mxexuse{39}

\begin{homework}Redo practice exercise 2 with the following data: replace the initial value with $315$ and the rate or ratio with $5$. Write it up in four lines: data, formula, calculation, conclusion.\end{homework}
\begin{solution}The solution depends on the exercise redone. Check above all the consistency of the model, the number of iterations, the rounding and the concluding sentence.\end{solution}

\begin{exitticket}[3 minutes]On a half-sheet: write the main rule of the session, work through a very short example and indicate your confidence level from 1 to 4. Key question(s) of the day: 1, 2 and 3.\end{exitticket}

\ifmxteacher
\begin{method}
\textbf{One-hour plan:} 7 min of fluency; 10 min of pooling and written record; 28 min of essential exercises; 5 min of correction; 3 min of exit ticket.\\
\textbf{Two-hour plan:} add 20 min of investigation, 20 min of formalisation, a second independent exercise and 10 min of comparative correction.\\
\textbf{Differentiation:} give struggling students a four-box sheet (data, model, calculation, conclusion). Ask fast students to produce a counterexample, a justification or a second method.
\end{method}
\fi
\mxsession{40}{\mxdate*{2027-03-11}}{B: 1 h -- C: 2 h}{Weighted trees}{Conditional probability}
    \begin{objectives}\begin{itemize}\item Construct and complete a tree.\item Read a conditional probability.\item Check the sums.\end{itemize}\end{objectives}
    \begin{fluency}[6 to 8 minutes, no calculator]\begin{enumerate}\item Write the multiplier for a $25\,\%$ increase. \mxblank{$1+\dfrac{25}{100}=1.25$}\item Calculate $0.8\times0.4$. \mxblank{$0.32$}\item If $P(A)=0.37$, calculate $P(\overline A)$. \mxblank{$0.63$}\item Simplify $\dfrac{18}{30}$. \mxblank{$\dfrac35$}\item What is the sum of the branches leaving a node? \mxblank{$1$}\end{enumerate}\end{fluency}
    \begin{prerequisites}Reread the last written record and be able to carry out the basic calculations used today. Before the investigation, have the keywords restated: construct and complete a tree and read a conditional probability.\end{prerequisites}

    \begin{activity}
    $35\,\%$ of orders are express. Among them, $4\,\%$ are late; among standard orders, $9\,\%$ are late. Construct and complete the tree.
    \begin{center}\begin{tikzpicture}[grow=right,level distance=30mm,sibling distance=10mm,every node/.style={font=\small},edge from parent/.style={draw,-{Stealth}}]
\node {$\Omega$} child {node {$A$} child {node {$B$}} child {node {$\overline B$}}} child {node {$\overline A$} child {node {$B$}} child {node {$\overline B$}}};
\end{tikzpicture}\end{center}
    \end{activity}
    \begin{solution}First branches: $0.35$ and $0.65$. Then $0.04$ and $0.96$; finally $0.09$ and $0.91$. Have students verbalise the piece of data that triggers the choice of method.\end{solution}

    \begin{summary}In a tree, the branches leaving a node sum to $1$. The second-level probabilities are conditional.\end{summary}

    \begin{example}
    Take the data of the investigation and write the solution in order: definition of the objects, property used, exact or approximate calculation, then a sentence giving the answer. The model answer appears in the teacher edition.
    \end{example}
    \begin{solution}First branches: $0.35$ and $0.65$. Then $0.04$ and $0.96$; finally $0.09$ and $0.91$.\end{solution}

    \begin{exercise}[essential]
    \begin{enumerate}\item We know $P(A)=0.40$, $P_A(B)=0.2$ and $P_{\overline A}(B)=0.4$. Construct the tree.\item Calculate $P(A\cap B)$ then $P(B)$.\item Calculate $P_B(A)$ and test independence.\end{enumerate}
    \end{exercise}
    \begin{solution}$P(A\cap B)=0.40\times0.2$. Then $P(B)=P(A)P_A(B)+(1-P(A))P_{\overline A}(B)$ and $P_B(A)=\dfrac{P(A\cap B)}{P(B)}$. Compare with $P(A)$.\end{solution}

    \mxfurther{Calao, course pp. 133--137; exercises p. 140, nos.~14--18}
\mxexuse{40}

    \begin{homework}Redo practice exercise 3 with the following data: replace the initial value with $320$ and the rate or ratio with $6$. Write it up in four lines: data, formula, calculation, conclusion.\end{homework}
    \begin{solution}The solution depends on the exercise redone. Check above all the consistency of the model, the number of iterations, the rounding and the concluding sentence.\end{solution}

    \begin{exitticket}[3 minutes]On a half-sheet: write the main rule of the session, work through a very short example and indicate your confidence level from 1 to 4. Key question(s) of the day: 1 and 2.\end{exitticket}

    \ifmxteacher
    \begin{method}
    \textbf{One-hour plan:} 7 min of fluency; 10 min of pooling and written record; 28 min of essential exercises; 5 min of correction; 3 min of exit ticket.\\
\textbf{Two-hour plan:} add 20 min of investigation, 20 min of formalisation, a second independent exercise and 10 min of comparative correction.\\
\textbf{Differentiation:} give struggling students a four-box sheet (data, model, calculation, conclusion). Ask fast students to produce a counterexample, a justification or a second method.
    \end{method}
    \fi
\mxsession{41}{\mxdate*{2027-03-16}}{B: 2 h -- C: 1 h}{Product of probabilities along a path}{Conditional probability}
    \begin{objectives}\begin{itemize}\item Calculate an intersection.\item Combine paths.\item Translate an event.\end{itemize}\end{objectives}
    \begin{fluency}[6 to 8 minutes, no calculator]\begin{enumerate}\item Simplify $3x+9-x+2$. \mxblank{$2x+11$}\item If $P(A)=0.37$, calculate $P(\overline A)$. \mxblank{$0.63$}\item Simplify $\dfrac{18}{30}$. \mxblank{$\dfrac35$}\item What is the sum of the branches leaving a node? \mxblank{$1$}\item Calculate $0.9\times0.4$. \mxblank{$0.36$}\end{enumerate}\end{fluency}
    \begin{prerequisites}Reread the last written record and be able to carry out the basic calculations used today. Before the investigation, have the keywords restated: calculate an intersection and combine paths.\end{prerequisites}

    \begin{activity}
    Using the tree from session 40, calculate the probability of an order being express and late, then the total probability of a delay. Interpret both results.
    \begin{center}\begin{tikzpicture}[grow=right,level distance=30mm,sibling distance=10mm,every node/.style={font=\small},edge from parent/.style={draw,-{Stealth}}]
\node {$\Omega$} child {node {$A$} child {node {$B$}} child {node {$\overline B$}}} child {node {$\overline A$} child {node {$B$}} child {node {$\overline B$}}};
\end{tikzpicture}\end{center}
    \end{activity}
    \begin{solution}$P(E\cap R)=0.35\times0.04=0.014$. $P(R)=0.014+0.65\times0.09=0.0725$. Have students verbalise the piece of data that triggers the choice of method.\end{solution}

    \begin{summary}We multiply along a path: $P(A\cap B)=P(A)P_A(B)$. We add the mutually exclusive paths that realise the event.\end{summary}

    \begin{example}
    Take the data of the investigation and write the solution in order: definition of the objects, property used, exact or approximate calculation, then a sentence giving the answer. The model answer appears in the teacher edition.
    \end{example}
    \begin{solution}$P(E\cap R)=0.35\times0.04=0.014$. $P(R)=0.014+0.65\times0.09=0.0725$.\end{solution}

    \begin{exercise}[essential]
    \begin{enumerate}\item We know $P(A)=0.45$, $P_A(B)=0.3$ and $P_{\overline A}(B)=0.5$. Construct the tree.\item Calculate $P(A\cap B)$ then $P(B)$.\item Calculate $P_B(A)$ and test independence.\end{enumerate}
    \end{exercise}
    \begin{solution}$P(A\cap B)=0.45\times0.3$. Then $P(B)=P(A)P_A(B)+(1-P(A))P_{\overline A}(B)$ and $P_B(A)=\dfrac{P(A\cap B)}{P(B)}$. Compare with $P(A)$.\end{solution}

    \mxfurther{Calao, p. 140, nos.~14--18}
\mxexuse{41}

    \begin{homework}Redo practice exercise 2 with the following data: replace the initial value with $325$ and the rate or ratio with $7$. Write it up in four lines: data, formula, calculation, conclusion.\end{homework}
    \begin{solution}The solution depends on the exercise redone. Check above all the consistency of the model, the number of iterations, the rounding and the concluding sentence.\end{solution}

    \begin{exitticket}[3 minutes]On a half-sheet: write the main rule of the session, work through a very short example and indicate your confidence level from 1 to 4. Key question(s) of the day: 1, 2 and 3.\end{exitticket}

    \ifmxteacher
    \begin{method}
    \textbf{One-hour plan:} 7 min of fluency; 10 min of pooling and written record; 28 min of essential exercises; 5 min of correction; 3 min of exit ticket.\\
\textbf{Two-hour plan:} add 20 min of investigation, 20 min of formalisation, a second independent exercise and 10 min of comparative correction.\\
\textbf{Differentiation:} give struggling students a four-box sheet (data, model, calculation, conclusion). Ask fast students to produce a counterexample, a justification or a second method.
    \end{method}
    \fi
\mxsession{42}{\mxdate*{2027-03-18}}{B: 1 h -- C: 2 h}{Independence}{Conditional probability}
\begin{objectives}\begin{itemize}\item Test independence.\item Compare probabilities.\item Justify a conclusion.\end{itemize}\end{objectives}
\begin{fluency}[6 to 8 minutes, no calculator]\begin{enumerate}\item Calculate $5\,\%$ of $60$. \mxblank{$3$}\item Simplify $\dfrac{18}{30}$. \mxblank{$\dfrac35$}\item What is the sum of the branches leaving a node? \mxblank{$1$}\item Calculate $0.3\times0.4$. \mxblank{$0.12$}\item If $P(A)=0.37$, calculate $P(\overline A)$. \mxblank{$0.63$}\end{enumerate}\end{fluency}
\begin{prerequisites}Reread the last written record and be able to carry out the basic calculations used today. Before the investigation, have the keywords restated: test independence and compare probabilities.\end{prerequisites}

\begin{activity}
In a population, $P(A)=0.4$, $P(B)=0.3$ and $P(A\cap B)=0.12$. Test independence in two ways. Then change $P(A\cap B)$ to $0.15$ and start again.

\end{activity}
\begin{solution}$0.4\times0.3=0.12$: independence. With $0.15$, the events are no longer independent. Have students verbalise the piece of data that triggers the choice of method.\end{solution}

\begin{summary}$A$ and $B$ are independent if $P(A\cap B)=P(A)P(B)$, equivalently if $P_A(B)=P(B)$ when $P(A)>0$.\end{summary}

\begin{example}
Take the data of the investigation and write the solution in order: definition of the objects, property used, exact or approximate calculation, then a sentence giving the answer. The model answer appears in the teacher edition.
\end{example}
\begin{solution}$0.4\times0.3=0.12$: independence. With $0.15$, the events are no longer independent.\end{solution}

\begin{exercise}[essential]
\begin{enumerate}\item We know $P(A)=0.20$, $P_A(B)=0.4$ and $P_{\overline A}(B)=0.6$. Construct the tree.\item Calculate $P(A\cap B)$ then $P(B)$.\item Calculate $P_B(A)$ and test independence.\end{enumerate}
\end{exercise}
\begin{solution}$P(A\cap B)=0.20\times0.4$. Then $P(B)=P(A)P_A(B)+(1-P(A))P_{\overline A}(B)$ and $P_B(A)=\dfrac{P(A\cap B)}{P(B)}$. Compare with $P(A)$.\end{solution}

\mxfurther{Calao, course pp. 134--137; exercises p. 144, nos.~33--35}
\mxexuse{42}

\begin{homework}Redo practice exercise 3 with the following data: replace the initial value with $330$ and the rate or ratio with $2$. Write it up in four lines: data, formula, calculation, conclusion.\end{homework}
\begin{solution}The solution depends on the exercise redone. Check above all the consistency of the model, the number of iterations, the rounding and the concluding sentence.\end{solution}

\begin{exitticket}[3 minutes]On a half-sheet: write the main rule of the session, work through a very short example and indicate your confidence level from 1 to 4. Key question(s) of the day: 1 and 2.\end{exitticket}

\ifmxteacher
\begin{method}
\textbf{One-hour plan:} 7 min of fluency; 10 min of pooling and written record; 28 min of essential exercises; 5 min of correction; 3 min of exit ticket.\\
\textbf{Two-hour plan:} add 20 min of investigation, 20 min of formalisation, a second independent exercise and 10 min of comparative correction.\\
\textbf{Differentiation:} give struggling students a four-box sheet (data, model, calculation, conclusion). Ask fast students to produce a counterexample, a justification or a second method.
\end{method}
\fi
\mxsession{43}{\mxdate*{2027-03-23}}{B: 2 h -- C: 1 h}{Partition and total probability}{Conditional probability}
    \begin{objectives}\begin{itemize}\item Recognise a partition.\item Apply total probability.\item Organise the calculation.\end{itemize}\end{objectives}
    \begin{fluency}[6 to 8 minutes, no calculator]\begin{enumerate}\item Write the multiplier for a $10\,\%$ increase. \mxblank{$1+\dfrac{10}{100}=1.10$}\item What is the sum of the branches leaving a node? \mxblank{$1$}\item Calculate $0.4\times0.4$. \mxblank{$0.16$}\item If $P(A)=0.37$, calculate $P(\overline A)$. \mxblank{$0.63$}\item Simplify $\dfrac{18}{30}$. \mxblank{$\dfrac35$}\end{enumerate}\end{fluency}
    \begin{prerequisites}Reread the last written record and be able to carry out the basic calculations used today. Before the investigation, have the keywords restated: recognise a partition and apply total probability.\end{prerequisites}

    \begin{activity}
    A factory has three lines A, B, C producing $20\,\%$, $50\,\%$ and $30\,\%$ of the items respectively. Their defect rates are $1\,\%$, $2\,\%$ and $4\,\%$. Calculate the overall defect rate.
    \begin{center}\begin{tikzpicture}[grow=right,level distance=30mm,sibling distance=10mm,every node/.style={font=\small},edge from parent/.style={draw,-{Stealth}}]
\node {$\Omega$} child {node {$A$} child {node {$B$}} child {node {$\overline B$}}} child {node {$\overline A$} child {node {$B$}} child {node {$\overline B$}}};
\end{tikzpicture}\end{center}
    \end{activity}
    \begin{solution}$P(D)=0.2\times0.01+0.5\times0.02+0.3\times0.04=0.024$, that is $2.4\,\%$. Have students verbalise the piece of data that triggers the choice of method.\end{solution}

    \begin{summary}If $(A_i)$ is a partition, $P(B)=\sum_i P(A_i)P_{A_i}(B)$. Each term corresponds to a path leading to $B$.\end{summary}

    \begin{example}
    Take the data of the investigation and write the solution in order: definition of the objects, property used, exact or approximate calculation, then a sentence giving the answer. The model answer appears in the teacher edition.
    \end{example}
    \begin{solution}$P(D)=0.2\times0.01+0.5\times0.02+0.3\times0.04=0.024$, that is $2.4\,\%$.\end{solution}

    \begin{exercise}[essential]
    \begin{enumerate}\item We know $P(A)=0.25$, $P_A(B)=0.5$ and $P_{\overline A}(B)=0.7$. Construct the tree.\item Calculate $P(A\cap B)$ then $P(B)$.\item Calculate $P_B(A)$ and test independence.\end{enumerate}
    \end{exercise}
    \begin{solution}$P(A\cap B)=0.25\times0.5$. Then $P(B)=P(A)P_A(B)+(1-P(A))P_{\overline A}(B)$ and $P_B(A)=\dfrac{P(A\cap B)}{P(B)}$. Compare with $P(A)$.\end{solution}

    \mxfurther{Calao, p. 138, nos.~1--5, then p. 142, nos.~22--25}
\mxexuse{43}

    \begin{homework}Redo practice exercise 2 with the following data: replace the initial value with $335$ and the rate or ratio with $3$. Write it up in four lines: data, formula, calculation, conclusion.\end{homework}
    \begin{solution}The solution depends on the exercise redone. Check above all the consistency of the model, the number of iterations, the rounding and the concluding sentence.\end{solution}

    \begin{exitticket}[3 minutes]On a half-sheet: write the main rule of the session, work through a very short example and indicate your confidence level from 1 to 4. Key question(s) of the day: 1, 2 and 3.\end{exitticket}

    \ifmxteacher
    \begin{method}
    \textbf{One-hour plan:} 7 min of fluency; 10 min of pooling and written record; 28 min of essential exercises; 5 min of correction; 3 min of exit ticket.\\
\textbf{Two-hour plan:} add 20 min of investigation, 20 min of formalisation, a second independent exercise and 10 min of comparative correction.\\
\textbf{Differentiation:} give struggling students a four-box sheet (data, model, calculation, conclusion). Ask fast students to produce a counterexample, a justification or a second method.
    \end{method}
    \fi
\mxsession{44}{\mxdate*{2027-03-25}}{B: 1 h -- C: 2 h}{Problems in context}{Conditional probability}
\begin{objectives}\begin{itemize}\item Choose a tree or a table.\item Solve a problem in context.\item Interpret.\end{itemize}\end{objectives}
\begin{fluency}[6 to 8 minutes, no calculator]\begin{enumerate}\item Simplify $3x+5-x+2$. \mxblank{$2x+7$}\item Calculate $0.5\times0.4$. \mxblank{$0.20$}\item If $P(A)=0.37$, calculate $P(\overline A)$. \mxblank{$0.63$}\item Simplify $\dfrac{18}{30}$. \mxblank{$\dfrac35$}\item What is the sum of the branches leaving a node? \mxblank{$1$}\end{enumerate}\end{fluency}
\begin{prerequisites}Reread the last written record and be able to carry out the basic calculations used today. Before the investigation, have the keywords restated: choose a tree or a table and solve a problem in context.\end{prerequisites}

\begin{activity}
Turn a table of counts (customers/subscription/cancellation) into a tree, then answer a question in each direction: probability of cancelling given the subscription; probability of having the subscription given the cancellation.

\end{activity}
\begin{solution}The first reading uses a row or a branch; the second requires dividing the intersection by the total number of cancellations. Have students verbalise the piece of data that triggers the choice of method.\end{solution}

\begin{summary}Choosing between a tree and a table: the tree highlights the order of conditioning; the table suits counts and reversed readings well.\end{summary}

\begin{example}
Take the data of the investigation and write the solution in order: definition of the objects, property used, exact or approximate calculation, then a sentence giving the answer. The model answer appears in the teacher edition.
\end{example}
\begin{solution}The first reading uses a row or a branch; the second requires dividing the intersection by the total number of cancellations.\end{solution}

\begin{exercise}[essential]
\begin{enumerate}\item We know $P(A)=0.30$, $P_A(B)=0.6$ and $P_{\overline A}(B)=0.8$. Construct the tree.\item Calculate $P(A\cap B)$ then $P(B)$.\item Calculate $P_B(A)$ and test independence.\end{enumerate}
\end{exercise}
\begin{solution}$P(A\cap B)=0.30\times0.6$. Then $P(B)=P(A)P_A(B)+(1-P(A))P_{\overline A}(B)$ and $P_B(A)=\dfrac{P(A\cap B)}{P(B)}$. Compare with $P(A)$.\end{solution}

\mxfurther{Calao, p. 142, nos.~22--25, then p. 144, nos.~30--32}
\mxexuse{44}

\begin{homework}Redo practice exercise 3 with the following data: replace the initial value with $340$ and the rate or ratio with $4$. Write it up in four lines: data, formula, calculation, conclusion.\end{homework}
\begin{solution}The solution depends on the exercise redone. Check above all the consistency of the model, the number of iterations, the rounding and the concluding sentence.\end{solution}

\begin{exitticket}[3 minutes]On a half-sheet: write the main rule of the session, work through a very short example and indicate your confidence level from 1 to 4. Key question(s) of the day: 1 and 2.\end{exitticket}

\ifmxteacher
\begin{method}
\textbf{One-hour plan:} 7 min of fluency; 10 min of pooling and written record; 28 min of essential exercises; 5 min of correction; 3 min of exit ticket.\\
\textbf{Two-hour plan:} add 20 min of investigation, 20 min of formalisation, a second independent exercise and 10 min of comparative correction.\\
\textbf{Differentiation:} give struggling students a four-box sheet (data, model, calculation, conclusion). Ask fast students to produce a counterexample, a justification or a second method.
\end{method}
\fi
\mxsession{45}{\mxdate*{2027-03-30}}{B: 2 h -- C: 1 h}{Mixed practice}{Conditional probability}
    \begin{objectives}\begin{itemize}\item Calculate a reverse probability.\item Chain total probability and conditioning.\item Check the result.\end{itemize}\end{objectives}
    \begin{fluency}[6 to 8 minutes, no calculator]\begin{enumerate}\item Calculate $20\,\%$ of $120$. \mxblank{$24$}\item If $P(A)=0.37$, calculate $P(\overline A)$. \mxblank{$0.63$}\item Simplify $\dfrac{18}{30}$. \mxblank{$\dfrac35$}\item What is the sum of the branches leaving a node? \mxblank{$1$}\item Calculate $0.6\times0.4$. \mxblank{$0.24$}\end{enumerate}\end{fluency}
    \begin{prerequisites}Reread the last written record and be able to carry out the basic calculations used today. Before the investigation, have the keywords restated: calculate a reverse probability and chain total probability and conditioning.\end{prerequisites}

    \begin{activity}
    In the situation of session 43, given that an item is defective, calculate the probability that it comes from line C. Explain why this result is not $30\,\%$.
    \begin{center}\begin{tikzpicture}[grow=right,level distance=30mm,sibling distance=10mm,every node/.style={font=\small},edge from parent/.style={draw,-{Stealth}}]
\node {$\Omega$} child {node {$A$} child {node {$B$}} child {node {$\overline B$}}} child {node {$\overline A$} child {node {$B$}} child {node {$\overline B$}}};
\end{tikzpicture}\end{center}
    \end{activity}
    \begin{solution}$P_D(C)=\dfrac{0.3\times0.04}{0.024}=0.5$. Conditioning changes the distribution because the defect rates differ. Have students verbalise the piece of data that triggers the choice of method.\end{solution}

    \begin{summary}For a reverse probability, first calculate the total probability of the conditioning event, then use the definition of conditional probability.\end{summary}

    \begin{example}
    Take the data of the investigation and write the solution in order: definition of the objects, property used, exact or approximate calculation, then a sentence giving the answer. The model answer appears in the teacher edition.
    \end{example}
    \begin{solution}$P_D(C)=\dfrac{0.3\times0.04}{0.024}=0.5$. Conditioning changes the distribution because the defect rates differ.\end{solution}

    \begin{exercise}[essential]
    \begin{enumerate}\item We know $P(A)=0.35$, $P_A(B)=0.2$ and $P_{\overline A}(B)=0.4$. Construct the tree.\item Calculate $P(A\cap B)$ then $P(B)$.\item Calculate $P_B(A)$ and test independence.\end{enumerate}
    \end{exercise}
    \begin{solution}$P(A\cap B)=0.35\times0.2$. Then $P(B)=P(A)P_A(B)+(1-P(A))P_{\overline A}(B)$ and $P_B(A)=\dfrac{P(A\cap B)}{P(B)}$. Compare with $P(A)$.\end{solution}

    \mxfurther{Calao, pp. 144--146}
\mxexuse{45}

    \begin{homework}Redo practice exercise 2 with the following data: replace the initial value with $345$ and the rate or ratio with $5$. Write it up in four lines: data, formula, calculation, conclusion.\end{homework}
    \begin{solution}The solution depends on the exercise redone. Check above all the consistency of the model, the number of iterations, the rounding and the concluding sentence.\end{solution}

    \begin{exitticket}[3 minutes]On a half-sheet: write the main rule of the session, work through a very short example and indicate your confidence level from 1 to 4. Key question(s) of the day: 1, 2 and 3.\end{exitticket}

    \ifmxteacher
    \begin{method}
    \textbf{One-hour plan:} 7 min of fluency; 10 min of pooling and written record; 28 min of essential exercises; 5 min of correction; 3 min of exit ticket.\\
\textbf{Two-hour plan:} add 20 min of investigation, 20 min of formalisation, a second independent exercise and 10 min of comparative correction.\\
\textbf{Differentiation:} give struggling students a four-box sheet (data, model, calculation, conclusion). Ask fast students to produce a counterexample, a justification or a second method.
    \end{method}
    \fi
\mxsession{46}{\mxdate*{2027-04-01}}{B: 1 h -- C: 2 h}{Review and assessment}{Conditional probability}
\begin{objectives}\begin{itemize}\item Use all the methods.\item Write up with the notation.\item Remedy mistakes.\end{itemize}\end{objectives}
\begin{fluency}[6 to 8 minutes, no calculator]\begin{enumerate}\item Write the multiplier for a $25\,\%$ increase. \mxblank{$1+\dfrac{25}{100}=1.25$}\item Simplify $\dfrac{18}{30}$. \mxblank{$\dfrac35$}\item What is the sum of the branches leaving a node? \mxblank{$1$}\item Calculate $0.7\times0.4$. \mxblank{$0.28$}\item If $P(A)=0.37$, calculate $P(\overline A)$. \mxblank{$0.63$}\end{enumerate}\end{fluency}
\begin{prerequisites}Reread the last written record and be able to carry out the basic calculations used today. Before the investigation, have the keywords restated: use all the methods and write up with the notation.\end{prerequisites}

\begin{activity}
Complete a four-station circuit: complete a tree; calculate an intersection; apply the total probability formula; test independence. Each group leaves a written justification for the next group.

\end{activity}
\begin{solution}The work produced must contain formulas and interpretations. Mistakes are sorted into: notation, complement, product, sum, reversed conditioning. Have students verbalise the piece of data that triggers the choice of method.\end{solution}

\begin{summary}Always name the events, complete all the branches, write the formula before calculating and interpret the final probability.\end{summary}

\begin{example}
Take the data of the investigation and write the solution in order: definition of the objects, property used, exact or approximate calculation, then a sentence giving the answer. The model answer appears in the teacher edition.
\end{example}
\begin{solution}The work produced must contain formulas and interpretations. Mistakes are sorted into: notation, complement, product, sum, reversed conditioning.\end{solution}

\begin{exercise}[essential]
\begin{enumerate}\item We know $P(A)=0.40$, $P_A(B)=0.3$ and $P_{\overline A}(B)=0.5$. Construct the tree.\item Calculate $P(A\cap B)$ then $P(B)$.\item Calculate $P_B(A)$ and test independence.\end{enumerate}
\end{exercise}
\begin{solution}$P(A\cap B)=0.40\times0.3$. Then $P(B)=P(A)P_A(B)+(1-P(A))P_{\overline A}(B)$ and $P_B(A)=\dfrac{P(A\cap B)}{P(B)}$. Compare with $P(A)$.\end{solution}

\mxfurther{Calao, pp. 146--147}
\mxexuse{46}

\begin{homework}Redo practice exercise 3 with the following data: replace the initial value with $350$ and the rate or ratio with $6$. Write it up in four lines: data, formula, calculation, conclusion.\end{homework}
\begin{solution}The solution depends on the exercise redone. Check above all the consistency of the model, the number of iterations, the rounding and the concluding sentence.\end{solution}

\begin{exitticket}[3 minutes]On a half-sheet: write the main rule of the session, work through a very short example and indicate your confidence level from 1 to 4. Key question(s) of the day: 1 and 2.\end{exitticket}

\ifmxteacher
\begin{method}
\textbf{One-hour plan:} 7 min of fluency; 10 min of pooling and written record; 28 min of essential exercises; 5 min of correction; 3 min of exit ticket.\\
\textbf{Two-hour plan:} add 20 min of investigation, 20 min of formalisation, a second independent exercise and 10 min of comparative correction.\\
\textbf{Differentiation:} give struggling students a four-box sheet (data, model, calculation, conclusion). Ask fast students to produce a counterexample, a justification or a second method.
\end{method}
\fi
